\chapter{Generalized Tonnetze and their topology}\label{ch:topology}

\section{From one trichord to a whole space}\label{sec:gentonnetz}
The classical Tonnetz is built from a single chord type: all transpositions and inversions of $\{0,4,7\}$. Nothing
prevents us from starting with any other trichord. Following Catanzaro \citep{catanzaro2011}, fix an equal
temperament $\Z_N$ and a partition $N=n_1+n_2+n_3$ into positive parts, and let
\[
\mathcal C(n_1,n_2,n_3)=\text{the chord complex whose triangles are all } T_t\text{ and } I_t\text{ images of }
\{0,\;n_1,\;n_1+n_2\}.
\]
Its vertices are the notes that occur, its edges are the note pairs that occur inside some chord, and its triangles
are the chords. The classical Tonnetz is $\mathcal C(3,4,5)$ in $\Z_{12}$. Related constructions include Cohn's
parsimonious trichords \citep{cohn1997}, Tymoczko's generalized Tonnetz of voice-leading lattices
\citep{tymoczko2012}, Bigo et al.'s chord-based complexes \citep{bigo2013}, and Yust's Tonnetze and ``Zeitnetze''
\citep{yust2020}; Jevti\'c and \v{Z}ivaljevi\'c prove that sufficiently generic higher-dimensional analogues
triangulate tori $T^{k-1}$ \citep{jevtic2021}.

Catanzaro's theorem \citep{catanzaro2011} classifies all these spaces for every $N$: the torus is the most common
outcome, especially as $N$ grows; the other possibilities are single 2-simplices, boundaries of tetrahedra, the
``harmonic strip'' (a cylinder or M\"obius band), and, for trichords containing a tritone, a ``circle of
tetrahedra boundaries'' related to Peck's Klein-bottle Tonnetz.

\section{All twelve cases for $N=12$, computed}
Up to transposition and inversion there are twelve trichord types in 12-TET. \Cref{tab:gentonnetz} lists the
complex of each, computed with \code{generalized\_tonnetz} and classified by the procedure of
Chapter~\ref{ch:prelims}: count components, check that every edge lies in at most two triangles, compute $\chi$
and orientability, then apply \cref{thm:classification}.

\begin{table}[t]\centering\small
\begin{tabular}{@{}llccccl@{}}\toprule
$(n_1,n_2,n_3)$ & Example chord & $V$ & $E$ & $F$ & $\chi$ & Space (Betti numbers over $\QQ$)\\\midrule
(1,1,10) & C--C$\sharp$--D cluster & 12 & 24 & 12 & 0 & cylinder (harmonic strip); (1,1,0)\\
(1,2,9)  & C--C$\sharp$--D$\sharp$ & 12 & 36 & 24 & 0 & torus; (1,2,1)\\
(1,3,8)  & C--C$\sharp$--E & 12 & 36 & 24 & 0 & torus; (1,2,1)\\
(1,4,7)  & C--C$\sharp$--F & 12 & 36 & 24 & 0 & torus; (1,2,1)\\
(1,5,6)  & C--C$\sharp$--F$\sharp$ & 12 & 30 & 24 & 6 & not a surface (edges in 4 triangles); (1,1,6)\\
(2,2,8)  & C--D--E & 12 & 24 & 12 & 0 & two cylinders (one per whole-tone scale); (2,2,0)\\
(2,3,7)  & C--D--F & 12 & 36 & 24 & 0 & torus; (1,2,1)\\
(2,4,6)  & C--D--F$\sharp$ & 12 & 30 & 24 & 6 & two components, not surfaces; (2,2,6)\\
(2,5,5)  & C--D--G & 12 & 24 & 12 & 0 & cylinder; (1,1,0)\\
(3,3,6)  & diminished C--E$\flat$--G$\flat$ & 12 & 18 & 12 & 6 & three tetrahedron boundaries (spheres); (3,0,3)\\
(3,4,5)  & major/minor triads & 12 & 36 & 24 & 0 & \textbf{torus}; (1,2,1)\\
(4,4,4)  & augmented C--E--G$\sharp$ & 12 & 12 & 4 & 4 & four disjoint triangles; (4,0,0)\\\bottomrule
\end{tabular}
\caption{Generalized Tonnetze in 12-TET. Betti numbers over $\Z_2$ coincide with those over $\QQ$ in every row.
Reproduced in \sheet{05\_generalized\_tonnetze.ipynb}; consistent with Catanzaro's classification.}
\label{tab:gentonnetz}
\end{table}

The table rewards careful reading.
\begin{itemize}
\item \emph{Six of the twelve types give tori}, including the classical Tonnetz. For a torus, each vertex must have
six distinct neighbours ($\pm n_1,\pm n_2,\pm n_3$ distinct), giving $E=36$ and $F=24$.
\item \emph{Symmetric chords collapse the space.} The augmented triad is fixed by $T_4$, so only four exist and
they are disjoint. The diminished triad lives inside a diminished-seventh collection of four notes whose four
3-subsets form the boundary of a tetrahedron, a sphere with $\chi=2$; three such collections give $\chi=6$.
\item \emph{Repeated gaps give strips.} When two gaps are equal (and no gap is a tritone or a third of the octave), the chord's neighbours coincide in pairs and the
complex is a long strip closing up on itself: a cylinder for $N=12$.
\item \emph{Tritones break the surface.} If a gap is 6, the tritone edge $\{x,x+6\}$ equals $\{x+6,x+12\}$ and is
shared by four triangles, so no neighbourhood looks like a disc. These are Catanzaro's circles of tetrahedra
boundaries.
\end{itemize}

\section{A parity phenomenon: cylinder or M\"obius band}
The chromatic cluster $\mathcal C(1,1,N-2)$ is a strip of triangles $\{x,x+1,x+2\}$ with the twisted edges of
consecutive triangles glued. Whether the strip closes with or without a half-twist depends on the parity of $N$.
Computing orientability for $5\le N\le15$ gives
\[
\mathcal C(1,1,N-2)\cong\begin{cases}\text{cylinder}, & N\text{ even},\\ \text{M\"obius band}, & N\text{ odd}.\end{cases}
\]
This is Catanzaro's harmonic strip ``in both its cylinder and M\"obius band variants''. The M\"obius case is the
first non-orientable space in this report; Chapter~\ref{ch:nonorient} meets closed non-orientable surfaces.

\paragraph{Beyond 12-TET.} In 19-tone equal temperament, looping over all 51 transposition classes of trichords
with \code{generalized\_tonnetz(a,b,c,n=19)} gives 42 tori and 9 M\"obius bands: since 19 is prime and odd, no chord
has transpositional symmetry and all harmonic strips twist. This illustrates the ``torus is generic'' part of
Catanzaro's theorem.

\begin{remark}[Why compute homology at all?]
Euler characteristic and orientability suffice for closed surfaces, but many musical complexes are not surfaces.
Betti numbers still describe them: $\mathcal C(1,5,6)$ has $b_2=6$ over $\QQ$, six independent enclosed
2-cycles, one per tritone pair. Topological signatures of this kind underlie the ``compliance'' measure of Bigo
et al.\ \citep{bigo2013} and persistent-homology fingerprints of musical pieces \citep{bergomi2016}.
\end{remark}

\begin{notebookbox}
\sheet{05\_generalized\_tonnetze.ipynb}: regenerate \cref{tab:gentonnetz} as a data frame, watch the harmonic strip
twist for odd $N$, and extend the table to 19-TET.
\end{notebookbox}

\begin{qa}
\Qu{How many trichord types are there in 12-TET up to transposition only, and up to transposition and inversion?}
\A{19 and 12. Inversion identifies a type with its mirror image, e.g.\ major (4,3,5) with minor (3,4,5).}
\Qu{Why does $\mathcal C(4,4,4)$ have only four triangles?}
\A{The augmented triad is invariant under $T_4$ and under inversion, so its $24$ images coincide in groups of six:
$24/6=4$ distinct triads, pairwise disjoint.}
\Qu{Show that each component of $\mathcal C(3,3,6)$ is a sphere.}
\A{A component is a diminished-seventh collection $\{x,x+3,x+6,x+9\}$; every 3-subset is a diminished triad, so
the four triangles form a tetrahedron boundary: $V-E+F=4-6+4=2$.}
\Qu{Why is $\mathcal C(1,5,6)$ not a surface?}
\A{Its tritone edges each lie in four triangles; around such an edge the space looks like four half-planes meeting
along a line, not like a plane.}
\Qu{For which $N$ is the chromatic strip a M\"obius band?}
\A{For odd $N$; for even $N$ it is a cylinder.}
\Qu{Name two spaces in \cref{tab:gentonnetz} with the same Euler characteristic but different topology.}
\A{The torus and the cylinder both have $\chi=0$; the torus is closed with $b_1=2$, the cylinder has boundary and
$b_1=1$. Euler characteristic alone does not classify spaces with boundary or with several components.}
\end{qa}
